\documentclass[11pt]{article}
\usepackage[utf8]{inputenc}
\usepackage{amsmath,amssymb}
\usepackage[margin=1in]{geometry}
\usepackage{hyperref}
\emergencystretch=3em

\title{Convolution algebra and the exponential Lipschitz estimate}
\author{Claude, with Daniel Murfet}
\date{2026-09-07}

\begin{document}
\maketitle
\sloppy

\begin{abstract}
First half of Astra~\#10's rank~5 (local Lipschitz dependence of the amplitude coefficients on the
double-series data). We record the linear algebra of the weighted convolution algebra
$(\ell^1_\rho(\mathbb N^2), *)$: bilinearity of \texttt{conv}, the right unit, closure of weighted
summability under sums and differences, and the triangle inequality for $\|\cdot\|_\rho$. On top of
this, the telescoping bound $\|a^{*n}-a'^{*n}\|_\rho\le nM^{n-1}\|a-a'\|_\rho$ for
$\|a\|_\rho,\|a'\|_\rho\le M$ and the Lipschitz estimate for the coefficient exponential
$\|E_a(t)-E_{a'}(t)\|_\rho\le |t|e^{|t|M}\|a-a'\|_\rho$ (\texttt{wnorm\_expCoeff\_sub\_le}).
Zero \texttt{sorry}/\texttt{axiom}.
\end{abstract}

\section{The things formalised}
{\small
\begin{verbatim}
theorem conv_sub_left (f g h) :
    conv (fun k => f k - g k) h = fun k => conv f h k - conv g h k
theorem conv_sub_right (f g h) :
    conv f (fun k => g k - h k) = fun k => conv f g k - conv f h k
theorem conv_delta_right (f) : conv f delta = f
theorem wsummable_add / wsummable_sub (ρ) (hρ : 0 ≤ ρ) (f g) (hf hg) : WSummable ρ (f ± g)
theorem wnorm_add_le (ρ) (hρ : 0 ≤ ρ) (f g) (hf hg) :
    wnorm ρ (fun k => f k + g k) ≤ wnorm ρ f + wnorm ρ g
theorem wnorm_convPow_sub_le (ρ) (hρ : 0 < ρ) (a a') (ha ha') (M) (hM hM') (n) :
    WSummable ρ (fun k => convPow a n k - convPow a' n k) ∧
    wnorm ρ (fun k => convPow a n k - convPow a' n k)
      ≤ if n = 0 then 0 else n * M ^ (n - 1) * wnorm ρ (fun k => a k - a' k)
theorem expDiffRow_hasSum (t M D) :
    HasSum (expDiffRow t M D) (|t| * Real.exp (|t| * M) * D)
theorem wnorm_expCoeff_sub_le (ρ) (hρ : 0 < ρ) (a a') (ha ha') (ha0 ha0') (M) (hM hM')
    (t) :
    WSummable ρ (fun k => expCoeff a t k - expCoeff a' t k) ∧
    wnorm ρ (fun k => expCoeff a t k - expCoeff a' t k)
      ≤ |t| * Real.exp (|t| * M) * wnorm ρ (fun k => a k - a' k)
\end{verbatim}
}

\section{Mathematics}
The convolution $(f*g)_k=\sum_{a\le k}f_a g_{k-a}$ is bilinear, so
\[
a^{*(n+1)}-a'^{*(n+1)}=(a-a')*a^{*n}+a'*(a^{*n}-a'^{*n}),
\]
and the submultiplicativity $\|f*g\|_\rho\le\|f\|_\rho\|g\|_\rho$ of unit~117 gives by induction
$\|a^{*n}-a'^{*n}\|_\rho\le nM^{n-1}\|a-a'\|_\rho$. No associativity or commutativity of $*$ is
needed for this: the identity uses only linearity in each slot. For the exponential
$E_a(t)_k=\sum_{n\le|k|}\frac{t^n}{n!}a^{*n}_k$ (with $a_{00}=0$, so the sum is really over all $n$),
the difference is $\sum_n\frac{t^n}{n!}(a^{*n}-a'^{*n})_k$, and the joint majorant
$g(n,k)=\frac{|t|^n}{n!}|a^{*n}_k-a'^{*n}_k|\rho^{|k|}$ has row sums
$\le\frac{|t|^n}{n!}nM^{n-1}D=|t|\frac{(|t|M)^{n-1}}{(n-1)!}D$ ($n\ge1$), summing to
$|t|e^{|t|M}D$ with $D=\|a-a'\|_\rho$. Tonelli on $\mathbb N\times\mathbb N^2$ then gives the
weighted bound. This is exactly the estimate $\|e^{a}-e^{a'}\|\le|t|e^{|t|M}\|a-a'\|$ one has in any
Banach algebra, specialised to the unital algebra of double power series with the $\rho$-weighted
$\ell^1$ norm.

\section{Paper-facing meaning}
The amplitude coefficients of §4.2 arise as $c(s)=y*E_{\bar x}(\beta s)$; the differentiability or
even continuity of the Taylor-tree coefficients in the data $(\xi,\eta)$ passes through a Lipschitz
estimate for $E$. Astra~\#10 ranked this fifth (after the parameter integration, measurability and
weighted-norm units, all landed). Unit~143 combines it with the scalar factor $e^{\beta s x_{00}}$
into the amplitude estimate proper.

\section{Lean notes}
\begin{itemize}
\item \texttt{HasSum.congr\_fun} takes the equation in the direction new $=$ old; the first attempt
with \texttt{.symm} was rejected with a type mismatch. Then \texttt{(hasSum\_nat\_add\_iff' 1).1}
with \texttt{simpa [h0]} (killing $\sum_{i<1}f_i=f_0=0$) transfers the shifted series to the
unshifted one; \texttt{convert} on a \texttt{HasSum} produced an unfillable
\texttt{AddCommMonoid} extensionality goal and should be avoided here.
\item The telescoping induction is cleanest stated for $n+1$ (\texttt{wnorm\_convPow\_succ\_sub\_le})
with the closed form $(n+1)M^n$; the $n=0$ case is the separate identity $\delta-\delta=0$, and the
general lemma carries an \texttt{if n = 0} guard so the row bound stays a single expression.
\item Only the \texttt{expCoeff\_eq\_tsum}-style rewrite (finite sum $=$ full \texttt{tsum} using
\texttt{convPow\_eq\_zero\_of\_lt}) is needed to compare with the row of the majorant; the rest is
the unit-118 majorant argument with \texttt{summable\_prod\_of\_nonneg} and
\texttt{Equiv.prodComm.tsum\_eq}.
\end{itemize}

\end{document}
